\documentclass[10pt,aspectratio=169]{beamer}
\usetheme{Madrid}
\usecolortheme{whale}
\usepackage[utf8]{inputenc}
\usepackage[T1]{fontenc}
\usepackage[spanish]{babel}
\usepackage{graphicx}
\usepackage{booktabs}
\usepackage{listings}

\definecolor{codebg}{rgb}{0.96,0.96,0.96}
\lstset{basicstyle=\ttfamily\scriptsize, backgroundcolor=\color{codebg},
        frame=single, breaklines=true, showstringspaces=false,
        literate={á}{{\'a}}1 {é}{{\'e}}1 {í}{{\'i}}1 {ó}{{\'o}}1 {ú}{{\'u}}1 {ñ}{{\~n}}1}

\title[Hito 2 - BD Literaria]{Proyecto Base de Datos I --- Hito 2}
\subtitle{Implementación, Optimización y Experimentación con Índices}
\author[Delgado, Guerrero, Melgarejo]{Sergio Delgado \and Antonio Guerrero \and Luciana Melgarejo}
\institute[UTEC]{Plataforma de catálogo literario}
\date{Ciclo 2026.1}

\begin{document}

\begin{frame}\titlepage\end{frame}

\begin{frame}{Agenda}
\tableofcontents
\end{frame}

\section{El proyecto}
\begin{frame}{El dominio}
\begin{columns}
\begin{column}{0.55\textwidth}
\textbf{Plataforma de catálogo literario:}
\begin{itemize}
  \item Entidad central: \texttt{Material} (libros, ensayos, revistas, poemas, audiolibros).
  \item Relacionada con autores, géneros, subgéneros, editoriales, premios.
  \item Interacciones de usuario: \texttt{Likes}, \texttt{Leer} (lecturas), \texttt{Resena}.
\end{itemize}
\vspace{0.3cm}
\textbf{25 tablas} en total: maestras, débiles, subtipos y tablas de relación.
\end{column}
\begin{column}{0.45\textwidth}
\footnotesize
\textbf{Del Hito 1 al Hito 2:}
\begin{itemize}
  \item Hito 1: modelo E-R y relacional.
  \item Hito 2: implementación física + experimento de índices.
\end{itemize}
\end{column}
\end{columns}
\end{frame}

\section{Implementación}
\begin{frame}{Modelamiento físico}
Tres mecanismos para las restricciones del Hito 1:
\begin{itemize}
  \item \textbf{Constraints} de columna/tabla: \texttt{PK}, \texttt{FK}, \texttt{UNIQUE}, \texttt{NOT NULL}, \texttt{CHECK}.
  \item \textbf{Discriminador} \texttt{Tipo} en \texttt{Material}: herencia total y disjunta.
  \item \textbf{Triggers} para reglas multi-tabla (4 funciones):
  \begin{itemize}\footnotesize
    \item Año de publicación $\geq$ fundación de la editorial.
    \item Fecha de reseña $\geq$ año del material.
    \item Todo material con $\geq$1 autor \emph{(trigger diferido)}.
    \item Todo género con $\geq$1 autor representante \emph{(trigger diferido)}.
  \end{itemize}
\end{itemize}
\vspace{0.2cm}
Más \textbf{5 vistas} de usuario que encapsulan los \texttt{JOIN} frecuentes.
\end{frame}

\begin{frame}{Cambios respecto al Hito 1}
\footnotesize
Cerramos reglas que el modelo relacional dejaba sin enforcement, sin reinterpretar el dominio:
\begin{itemize}
  \item \textbf{[A]} Discriminador \texttt{Tipo}; \textbf{[B][C]} triggers de participación total.
  \item \textbf{[D]} \texttt{Fecha} en PK de \texttt{Leer} (permite relecturas).
  \item \textbf{[I]} Nueva tabla \texttt{PerteneceSubGenero} (consulta "por subgénero" que el Hito 1 pedía pero el modelo dejaba imposible).
  \item \textbf{[J]} Reescritura de Q3 para evitar un producto cartesiano.
  \item \textbf{[K]} Quitamos índices redundantes con la PK.
\end{itemize}
\centering (11 cambios documentados en \texttt{docs/cambios\_respecto\_hito1.md})
\end{frame}

\begin{frame}{Datos: 4 escenarios reproducibles}
Generador propio (Python + Faker), carga por \texttt{COPY}, semilla fija.
\vspace{0.3cm}
\centering
\begin{tabular}{lrrr}
\toprule
\textbf{Base} & \textbf{Materiales} & \textbf{Likes/Lecturas/Reseñas} & \textbf{Total filas} \\
\midrule
1k   &     300 &     1\,000 c/u & $\sim$5 mil \\
10k  &   2\,000 &    10\,000 c/u & $\sim$46 mil \\
100k &  20\,000 &   100\,000 c/u & $\sim$484 mil \\
1m   & 100\,000 & 1\,000\,000 c/u & \textbf{$\sim$3.86 millones} \\
\bottomrule
\end{tabular}
\end{frame}

\section{Experimentación}
\begin{frame}{Las 5 consultas y la metodología}
\footnotesize
\begin{tabular}{lll}
\toprule
\textbf{ID} & \textbf{Qué responde} & \textbf{Patrón} \\
\midrule
Q1 & Materiales de un autor & igualdad selectiva \\
Q2 & Materiales de un género & igualdad poco selectiva \\
Q3 & Materiales más populares & agregación total \\
Q4 & Historial de un usuario & igualdad + orden \\
Q5 & Materiales en un \textbf{rango} de años & \textbf{rango (BETWEEN)} \\
\bottomrule
\end{tabular}
\vspace{0.3cm}

\textbf{Metodología:} \texttt{EXPLAIN (ANALYZE, BUFFERS)}, 6 corridas (se descarta la fría, mediana de 5). \texttt{DISCARD ALL} entre corridas. Escenario "sin índices" con \texttt{05\_drop\_indexes.sql}. Se mide con y sin índices en las 4 bases.
\end{frame}

\begin{frame}{Resultado global (base 1M)}
\begin{columns}
\begin{column}{0.5\textwidth}
\centering
\includegraphics[width=\textwidth]{img/speedup_1m.png}
\end{column}
\begin{column}{0.5\textwidth}
\footnotesize
\begin{tabular}{lrr}
\toprule
& \textbf{Sin} & \textbf{Con} \\
\midrule
Q1 (autor)     &    7.19 &  0.08 \\
Q5 (rango)     &    9.64 &  0.19 \\
Q4 (historial) &  131.05 &  8.00 \\
Q2 (género)    &   34.01 & 22.41 \\
Q3 (populares) & 2335.55 & 2345.02 \\
\bottomrule
\end{tabular}
\vspace{0.2cm}

\textbf{3 ganan mucho, 1 poco, 1 nada.}
\end{column}
\end{columns}
\end{frame}

\begin{frame}{Q5 --- La consulta por rango (caso estrella)}
\begin{columns}
\begin{column}{0.5\textwidth}
\includegraphics[width=\textwidth]{img/Q5_por_rango_anio.png}
\end{column}
\begin{column}{0.5\textwidth}
\footnotesize
\texttt{WHERE Anio\_Publicacion BETWEEN 1900 AND 1905}
\vspace{0.2cm}

\textbf{Sin índice:} \texttt{Seq Scan} de toda la tabla. Crece con el tamaño.

\textbf{Con índice:} \texttt{Bitmap Index Scan}, toca solo el tramo del rango.

\vspace{0.2cm}
$\Rightarrow$ \textbf{52× más rápido} en 1M, y casi plano al escalar.
\end{column}
\end{columns}
\end{frame}

\begin{frame}{Q3 --- Cuando el índice NO ayuda}
\begin{columns}
\begin{column}{0.5\textwidth}
\includegraphics[width=\textwidth]{img/Q3_populares.png}
\end{column}
\begin{column}{0.5\textwidth}
\footnotesize
Agregación \textbf{total} de likes, lecturas y reseñas: no hay filtro selectivo.
\vspace{0.2cm}

El motor debe leer las tablas completas con o sin índice.
\vspace{0.2cm}

\textbf{2335 ms vs 2345 ms} $\to$ 1.0×. Un índice solo ayuda si permite \emph{evitar} leer filas.
\end{column}
\end{columns}
\end{frame}

\begin{frame}{El índice no es gratis}
\begin{columns}
\begin{column}{0.5\textwidth}
\includegraphics[width=\textwidth]{img/overhead_escritura.png}
\end{column}
\begin{column}{0.5\textwidth}
\footnotesize
\textbf{Escrituras:} insertar 1M filas con 2 índices que mantener = \textbf{7.6× más lento} (557 ms $\to$ 4228 ms).
\vspace{0.3cm}

\textbf{Lectura no selectiva forzada:} usar el índice cuando se lee $\sim$100\% de la tabla = \textbf{2.4× más lento} que el \texttt{Seq Scan} (16.6 ms $\to$ 40.5 ms).
\vspace{0.3cm}

Por eso el planner lo evita: es la decisión correcta.
\end{column}
\end{columns}
\end{frame}

\section{Conclusiones}
\begin{frame}{Conclusiones}
\begin{itemize}
  \item Implementación completa del modelo del Hito 1: tablas, restricciones, triggers y vistas.
  \item Cuatro volúmenes reproducibles, hasta $\sim$3.86M de filas.
  \item El experimento confirma los tres comportamientos del índice: gran beneficio (hasta 91×), moderado (1.5×) y nulo (1.0×); más su costo en escritura (7.6×).
  \item \textbf{Indexar es una decisión por consulta y por selectividad}, no una receta universal.
  \item Por encima del millón: la arquitectura cliente-servidor basta para consultas selectivas bien indexadas; las agregaciones totales son el límite.
\end{itemize}
\end{frame}

\begin{frame}
\centering
\Large ¿Preguntas?
\vspace{0.5cm}

\footnotesize \url{https://github.com/SefedeamU/Proyecto-BD-Hito-2}
\end{frame}

\end{document}
